\section{Related Work}\label{sec:related}

Our work connects three research threads: self-consistency and test-time scaling methods developed for reasoning tasks, confidence estimation for structured information extraction, and broader uncertainty quantification for language models.

\subsection{Self-Consistency and Test-Time Scaling}

\paragraph{Self-consistency for discrete and free-form tasks.}
Self-consistency decoding \cite{wang2023selfconsistency} improves chain-of-thought reasoning by sampling multiple rationales and selecting the most frequent answer via majority vote.
This paradigm assumes a discrete answer space where outputs can be grouped by equivalence and counted---an assumption that holds for arithmetic, commonsense QA, and similar tasks but breaks for set-valued structured outputs.
\citet{chen2024usc} extend self-consistency to free-form text generation by replacing exact matching with LLM-based semantic comparison, while SelfCheckGPT \cite{manakul2023selfcheckgpt} leverages cross-sample consistency to detect hallucinations in open-ended generation.
\citet{thirukovalluru2024atomic} decompose long-form outputs into atomic sub-claims and measure per-claim consistency, moving beyond whole-output comparison.
These methods operate on text-level equivalence or entailment; our work instead compares \textit{structured sets} of entities and relations, requiring set-theoretic similarity measures that account for partial overlap, type constraints, and span-level matching.

\paragraph{Test-time scaling and sample selection.}
A growing body of work studies how to allocate additional compute at inference time.
\citet{snell2024scaling} characterize the scaling laws governing test-time compute, showing that repeated sampling with effective selection can substitute for model scaling.
\citet{brown2024monkeys} demonstrate that best-of-$N$ sampling with large $N$ yields substantial gains across diverse tasks.
Several recent methods aim to improve selection without trained reward models.
CISC \cite{taubenfeld2025cisc} combines confidence scores with majority voting, weighting each sample's contribution to the consensus.
Self-Certainty \cite{kang2025selfcertainty} uses KL divergence between a sample's token distribution and the prior to select high-quality outputs, demonstrating that reward-free selection is viable for mathematical reasoning.
ReASC \cite{kim2026reasc} adaptively adjusts the number of samples based on estimated difficulty, while \citet{wang2025rankedvoting} replace plurality voting with ranked aggregation, and \citet{oh2025latentsc} perform consistency estimation in latent space rather than output space.

All of these methods target reasoning and QA tasks with discrete or near-discrete answer spaces.
Our work is, to our knowledge, the first to extend self-consistency to set-valued structured outputs in IE, and the first to identify the correlation-selection gap: strong quality estimation ($\rho{=}0.50$) coexists with selection failure (no gain over greedy on NER).
In direct dialogue with \citet{kang2025selfcertainty}, who show that reward-free selection succeeds for reasoning, our results demonstrate that this success does not transfer to structured extraction, where the combinatorial output space and instance-level nature of consistency signals preclude effective per-sample discrimination.

\subsection{Consistency and Confidence for Structured IE}

The most directly relevant prior work is \citet{nedellec2026consistency}, who study consistency--accuracy correlations in hard-prompted LLMs for NER and relation extraction on plant-health data.
Their findings corroborate our quality-estimation results: higher inter-sample consistency predicts higher extraction accuracy.
However, three key limitations motivate our work.
First, they rely on surface-level string matching and explicitly acknowledge that their approach ``overlooks semantically consistent outputs''---precisely the gap our Soft Jaccard signal addresses through type-aware bipartite matching with span overlap.
Second, they do not study best-of-$N$ selection, leaving the quality-estimation-to-selection transfer question unexamined.
Third, their evaluation covers a single domain with hard-prompted (temperature-zero) generation and no comparison across signal families.
Our work extends their analysis to multi-domain benchmarks, temperature-based sampling, and a systematic five-signal comparison that reveals when structural signals outperform surface alternatives and vice versa.

CoCoA \cite{vashurin2025cocoa} proposes a minimum Bayes risk framework that fuses model confidence with inter-sample consistency for uncertainty quantification on QA, summarization, and translation.
Their confidence-consistency ensemble parallels our analysis of signal combinations, but they do not address structured extraction, where the unit of comparison shifts from text sequences to entity and relation sets.

Our evaluation setting---LLMs performing IE as a generation task---builds on a line of work establishing generative approaches to entity and relation extraction.
UniNER \cite{zhou2024uniner} and GoLLIE \cite{sainz2024gollie} train universal IE models via instruction tuning; \citet{wadhwa2023revisiting} and InstructUIE \cite{wang2023instructuie} demonstrate that instruction-following LLMs can perform extraction zero-shot or with minimal supervision; GenIE \cite{josifoski2022genie} and CodeIE \cite{li2023codeie} frame IE as constrained sequence generation.
We adopt constrained decoding \cite{willard2023guided} following this paradigm and study confidence estimation as an orthogonal dimension: given that LLMs generate structured outputs, how reliably can we estimate extraction quality from multiple samples?

\subsection{Uncertainty Estimation and Calibration}

\paragraph{Confidence and uncertainty for LLMs.}
Several surveys chart the landscape of confidence estimation \cite{geng2024confidence}, selective abstention \cite{wen2025know}, uncertainty quantification \cite{xia2025survey}, and their intersection with calibration \cite{hobelsberger2025systematic}.
Our work bridges confidence estimation and selective prediction: we use confidence signals not only to rank instances by reliability but also to attempt best-of-$N$ selection, revealing that these two objectives impose different requirements on the underlying signal.

Semantic uncertainty \cite{kuhn2023semantic} clusters generations by meaning rather than surface form and estimates uncertainty from cluster statistics.
\citet{kossen2024semantic} distill semantic entropy into lightweight probes, and P(True) \cite{kadavath2022language} directly queries the model for its confidence.
These approaches measure uncertainty over the model's belief distribution; our structural signals instead measure agreement over output \textit{structure}, comparing extracted entity sets rather than semantic content of free-form text.

\paragraph{Calibration and ensembles.}
Modern neural networks are often miscalibrated \cite{guo2017calibration}, and calibration techniques for pre-trained language models \cite{desai2020calibration} aim to align predicted probabilities with empirical accuracy.
Our reliability diagrams assess whether confidence signals are calibrated for IE, extending calibration analysis beyond classification to structured prediction.
Deep ensembles \cite{lakshminarayanan2017ensembles} and MC-dropout \cite{gal2016dropout} obtain uncertainty from model variation, requiring multiple forward passes through different model parameters.
Our multi-sample approach is conceptually related---we also use variation across stochastic outputs---but requires no model modification or access to internal parameters beyond token log-probabilities, operating entirely post-hoc on generated outputs.

\paragraph{Uncertainty for token-level and span-level NER.}
Prior work on NER uncertainty operates within discriminative, pre-LLM paradigms: \citet{xiao2019uncertainty} study Bayesian inference for sequence labeling, \citet{shelmanov2021active} leverage uncertainty for active learning in NER, and \citet{zhang2023ener} estimate entity-level uncertainty from span probabilities.
These methods assume access to per-token classification logits from a discriminative model.
In the generative LLM setting we study, entity boundaries are implicit in the decoded text and uncertainty must be estimated from the generation process itself.

\paragraph{Related paradigms.}
Conformal prediction \cite{fisch2021conformal, angelopoulos2021conformal} provides distribution-free coverage guarantees but targets prediction sets rather than quality ranking.
Trained verifiers \cite{cobbe2021verifiers, lightman2024verify} and reward models \cite{wang2024mathshepherd} learn to score outputs from labeled data; we focus on unsupervised signals that require no additional training.
Selective prediction \cite{geifman2017selective} and confidence-based IE filtering \cite{culotta2004confidence} establish the broader framework of abstaining when uncertain; our contribution is characterizing which confidence signals support this objective for generative IE and where the quality-estimation-to-selection transfer breaks down.
